\documentclass[aps,prb,twocolumn,superscriptaddress]{revtex4-2}
\usepackage{amsmath,amssymb,graphicx,hyperref}
\begin{document}
\title{Replication Report: Altermagnetic Spin Splitting in $\alpha$-MnTe\\
\small A from-scratch tight-binding surrogate of Sandratskii, Carva \& Silkin (2025), arXiv:2501.11327}
\author{REPLICATE-PROJECT (automated)}
\date{\today}
\begin{abstract}
We replicate the electronic-structure headline of Sandratskii \textit{et al.}
(arXiv:2501.11327v3): $\alpha$-MnTe is an altermagnet whose collinear
antiparallel Mn sublattices produce \emph{zero} net moment yet a nonzero,
momentum-dependent spin splitting $\Delta E(\mathbf{k})=\varepsilon_\uparrow(\mathbf{k})-\varepsilon_\downarrow(\mathbf{k})$
of the electron states. Rather than reproduce the full DFT magnon method, we
build a minimal from-scratch tight-binding altermagnet in which the two Mn
sublattices are related by a spin-space-group operation (real-space $C_{6z}$
rotation $+$ spin flip). The surrogate reproduces all five qualitative
predictions of the paper's electronic headline: (P1) exact moment compensation,
(P2) a spin-degenerate nodal line along $(0,0,k_z)$, (P3) nonzero splitting at
general $\mathbf{k}$, (P4) a sign-changing multi-lobe (wave-like) form factor
with vanishing Brillouin-zone average, and (P5) a well-defined $k_z$ parity.
An isotropic (ferromagnet-like) null model gives identically zero splitting,
confirming that the effect is symmetry-enforced. \textbf{Verdict: PARTIAL} ---
the electronic altermagnetism is reproduced; the surrogate's angular form factor
is $d$-wave-like (4 nodes) rather than MnTe's true $g$-wave (8 nodes), and the
DFT magnon method itself is not reproduced.
\end{abstract}
\maketitle

\section{Target claim}
The paper's abstract makes two coupled claims: (i) a new direct-DFT method for
adiabatic magnons in collinear magnets via a magnetization-change constraint plus
a generalized-translation symmetry constraint; and (ii) that in $\alpha$-MnTe the
altermagnetic \emph{electronic} spin splitting in $\mathbf{k}$-space and the
magnon \emph{chirality} splitting in $\mathbf{q}$-space share an identical
reciprocal-space pattern. The single most testable, cheaply reproducible headline
is the electronic one: a zero-net-moment collinear magnet with a symmetry-enforced,
sign-changing, momentum-dependent spin splitting.

\section{Method (surrogate)}
We do not reproduce the DFT. We build a two-sublattice hexagonal tight-binding
model for $\alpha$-MnTe (NiAs-type, $a=4.15$~\AA, $c=6.71$~\AA). Sublattices
$A,B$ carry antiparallel moments $\pm J$. The altermagnetic defining relation is
imposed exactly: $\varepsilon_\downarrow(\mathbf{k})=\varepsilon_\uparrow(R^{-1}\mathbf{k})$
with $R=C_{6z}$ the sublattice-connecting rotation, so
$\Delta E(\mathbf{k})=\varepsilon_\uparrow(\mathbf{k})-\varepsilon_\uparrow(R^{-1}\mathbf{k})$.
The up band $\varepsilon_\uparrow$ combines a $C_6$-symmetric isotropic hopping
(which cancels in $\Delta E$) with an anisotropic ``altermagnetic'' hopping
$\propto(2\cos k\!\cdot\!a_1-\cos k\!\cdot\!a_2-\cos k\!\cdot\!a_3)$ modulated
along $k_z$ by the $c/2$ screw stacking.

\section{Results}
All measured, none asserted (see \texttt{sandratskii2025\_result.json}):
\begin{itemize}
\item \textbf{P1} net moment $=0.0\,\mu_B$ (exact compensation).
\item \textbf{P2} $\max|\Delta E|$ on $(0,0,k_z)=0.0$: spin-degenerate nodal line.
\item \textbf{P3} $\max|\Delta E|$ on $(0.1,0.2,k_z)=0.072$ (model units): nonzero at general $\mathbf{k}$.
\item \textbf{P4} BZ-average $\Delta E / \max|\Delta E| = 0.020$ (sign-changing);
positive/negative area fractions $=0.500/0.500$; 4 angular sign changes at $|k|=0.3$.
\item \textbf{P5} $k_z$ parity well-defined (even, residual $0.0$).
\item \textbf{Null test} isotropic model $\max|\Delta E|=3\times10^{-15}$: no
splitting without the anisotropic altermagnetic term.
\end{itemize}
Predictions passed: 6/6 qualitative checks.

\section{Assessment}
The surrogate confirms the \emph{qualitative} electronic altermagnetism of the
paper's headline: zero net moment coexisting with a symmetry-enforced,
sign-changing, momentum-dependent spin splitting that is nodal on the
high-symmetry axis and averages to zero over the BZ. Two honest gaps remain:
(1) the surrogate's form factor is $d$-wave-like (4 nodes) whereas MnTe's true
symmetry gives a $g$-wave (8-node) pattern --- a form-factor mismatch, not a
qualitative failure; and (2) the paper's actual contribution (the direct-DFT
magnon method and the magnon--electron pattern equivalence) is beyond a
$<8$-minute model surrogate and is not reproduced.
\end{document}
